%% fig_xphi.tex
%% (x, phi) 占有マップ: 各位置 x（1周期に畳んだもの）での配向条件分布 P(phi|x)。
%% m=6, beta pE=3 で、delta=1/3（配向が電場方向に揃う）と delta=10（軸方向へ傾く）を比較。
%% 破線は最狭窄部 x~0.81 と最広部 x~0.19。点線は phi=pi/2（電場方向）。
\begin{tikzpicture}
  \begin{groupplot}[
      group style={group size=1 by 2, vertical sep=0.9cm,
        xticklabels at=edge bottom},
      width=8.4cm, height=3.6cm, scale only axis,
      view={0}{90},
      tick align=outside,
      tick label style={font=\footnotesize}, label style={font=\small},
      title style={font=\small},
      xmin=0, xmax=1, ymin=0, ymax=6.2832,
      ylabel={$\phi$},
      ytick={0,1.5708,3.1416,4.7124,6.2832},
      yticklabels={$0$,$\tfrac{\pi}{2}$,$\pi$,$\tfrac{3\pi}{2}$,$2\pi$},
      enlargelimits=false, grid=none,
      colormap/viridis,
      colorbar, colorbar style={font=\footnotesize, width=0.22cm,
        ylabel={$P(\phi\,|\,x)$}},
      point meta min=0, point meta max=1.2,
    ]
    \nextgroupplot[title={$\delta=|\Delta\alpha|E/p=1/3$（揃う）}]
      \addplot3[surf, shader=interp]
        table[x=x, y=phi, z=p]{\figdata/xphi_m6_b3_d0333333.dat};
      \draw[white, densely dashed, line width=0.7pt]
        (axis cs:0.81,0) -- (axis cs:0.81,6.2832);
      \draw[white, dotted, line width=0.7pt]
        (axis cs:0,1.5708) -- (axis cs:1,1.5708);
    \nextgroupplot[title={$\delta=10$（傾く）},
      xlabel={$x$ (1 周期に畳んだ位置)}]
      \addplot3[surf, shader=interp]
        table[x=x, y=phi, z=p]{\figdata/xphi_m6_b3_d10.dat};
      \draw[white, densely dashed, line width=0.7pt]
        (axis cs:0.81,0) -- (axis cs:0.81,6.2832);
      \draw[white, dotted, line width=0.7pt]
        (axis cs:0,1.5708) -- (axis cs:1,1.5708);
  \end{groupplot}
\end{tikzpicture}
